% LaTeX テンプレート — 授業資料用 (LuaLaTeX を想定、両面: 奇数ページにヘッダー)
\documentclass[11pt,twoside]{article}

% LuaLaTeX 向け設定: 日本語対応に luatexja-fontspec と fontspec を使用
\usepackage{fontspec}
\usepackage{luatexja-fontspec}

% 必要に応じて日本語フォントを明示（環境に応じて変更してください）
% 例: \setmainjfont{Noto Serif CJK JP}

% 単位表示の統一: siunitx を使用
\usepackage{siunitx}
\sisetup{detect-mode=true, per-mode=symbol}

% その他のパッケージ
\usepackage{tikz}
\usepackage{fancybox}
\usepackage{ascmac}
\usepackage{amsmath}
\usepackage{geometry}
\usepackage{fancyhdr}
\usepackage{lipsum} % サンプル用ダミーテキスト
\usepackage{ulem}
\usepackage{enumitem}
\usepackage{amssymb}
\usepackage{dashrule}
\usepackage{graphicx}

%Tikz
\usepackage{physmotion}

% 使い方: 以下を編集してください
\newcommand{\SubjectTitle}{斜方投射} %テーマ
\newcommand{\ClassRoomInfo}{\underline{\hspace{0.5cm}}年\underline{\hspace{0.5cm}}組　氏名\underline{\hspace{3cm}}}

% 余白設定
\geometry{top=20mm,bottom=25mm,left=20mm,right=20mm,heightrounded}

% ヘッダー設定: 奇数ページに表示（左: 科目と内容、右: 学年組欄）
\setlength{\headheight}{14pt}
\setlength{\headsep}{36pt}
\pagestyle{fancy}
\fancyhf{} % 既定のヘッダー/フッターをクリア
\fancyhead[LO]{\small 物理基礎「\SubjectTitle」} %Odd page Left headerを設定
\fancyhead[RO]{\small \ClassRoomInfo} % Odd page Right headerを設定
\renewcommand{\headrulewidth}{0.4pt}
\renewcommand{\footrulewidth}{0pt}
% フッター中央にページ番号を表示
\fancyfoot[C]{\small\thepage}

% 偶数ページはヘッダーを空にして、本文が上から始まるようにする
\fancyhead[LE,CE,RE]{}

% ドキュメント開始
\begin{document}
	% 斜方投射
	\section{斜方投射}
		\begin{minipage}[t]{0.44\linewidth}
			\obliqueprojectilediagram{v_0}{\theta}{v_{0x}}{v_{0y}}{g}{t}\\\\
			\begin{itembox}[l]{補足}
				I. 最高点について\\
				◎最高点では$v_y=0$(上図を参照)\\
				右の(ii)より
				\[
					v_y=v_0\sin\theta -gt
				\]
				これがゼロになるとき最高点なので
				\[
					t_{\text{最高点}}=\frac{v_0\sin\theta}{g}
				\]
				II. ボールの軌道の導出\\
				(i)-③式より
				\[
					t=\frac{x}{v_0\cos\theta}
				\]
				これを(ii)-③式に入れて計算することで、
				\[
					y=-\frac{g}{2v_0^2\cos^2\theta}x^2+(\tan\theta)x\quad\text{：二次関数}
				\]
			\end{itembox}\\\\
			\begin{itembox}[l]{<発展>最も遠くまで飛ばせる角度は？}
				今回学習したことに加えて、三角関数の知識を用いると、斜方投射で一番遠くまで飛ばすことのできる角度が計算できる。\\
				方針は、$x=v_{0}\cos\theta ×t$に着地する時間$t$を代入して、$x$を$\theta$の関数として表すことを目指す。\\
				そうした後、$x$が最大になる角度$\theta$を調べるともっとも遠くまで飛ばせる角度がわかる。

			\end{itembox}
			% 　\\
			% \begin{itembox}[l]{<参考>$v_0$,$v_x$,$v_{0x}$}
			% 	物理の問題では、変数の右下にラベルを\\張ることが多々ある。\\
			% 	今回出てきた$v_0$、$v_x$、$v_{0x}$もその例である。\\
			% 	以下にそれぞれの意味をまとめる。
			% 	\begin{center}
			% 		\begin{tabular}{ll}
			% 			$v_0$    & ：初速度の大きさ \\
			% 			$v_x$    & ：$x$軸方向の速度 \\
			% 			$v_{0x}$ & ：$x$軸方向の初速度 \\
			% 		\end{tabular}
			% 	\end{center}
			% 	　
			% \end{itembox}
		\end{minipage}
		\hfill	
		\begin{minipage}[t]{0.52\linewidth}
			\ovalbox{①斜方投射とは？}\\
			左図のような重力がかかる運動のうち、初速度$v_0$が任意の方向を向いている物体の運動。\\\\
			\ovalbox{②今までとの違い}\\
			$v_0$が$x$成分と$y$成分の両方の値を持ちうる。\\\\
			\ovalbox{③セットアップ}
				\begin{itemize}[leftmargin=2em,itemsep=0pt,topsep=2pt]
					\item[(i)]
					軸の向き：
					水平右向きを $x$ 軸の正、
					鉛直上向きを $y$ 軸の正とする。
	
					\item[(ii)]
					$t=0$ のとき、
					物体は原点 $(x,y)=(0,0)$ にある。
	
					\item[(iii)]
					重力加速度の大きさを $g$、
					初速度の大きさを $v_0$、仰角を $\theta$ とする。
	
					\item[(iv)]
					初速度の補足：各方向の初速度は以下のようになる。\\
					$v_{0x} = v_0\cos\theta$（$x$軸方向の初速度）\\
					$v_{0y} = v_0\sin\theta$（$y$軸方向の初速度）\\
				\end{itemize}
			\ovalbox{④$t$秒後の物体の位置と速度}　\uuline{\textbf{大事！！}}\\
			\textbf{point!：}$x$と$y$を分けて考える\\
			\textbf{point!：}$y$軸方向のみ加速している
			\begin{itemize}[leftmargin=2em,itemsep=0pt,topsep=2pt]
				\item[(i)]
				$x$軸方向の運動：\\
				\(
				\left(
				\begin{array}{ll}
					\text{①} \quad a_x = & \\
					\text{②} \quad v_x = & \\
					\text{③} \quad x = & \\
				\end{array}
				\right.
				\)\\
				\to $x$軸方向の運動は\underline{\hspace{3.5cm}}と一致！\\
				\item[(ii)]
				$y$軸方向の運動：\\
				\(
				\left(
				\begin{array}{ll}
					\text{①} \quad a_y = & \\
					\text{②} \quad v_y = & \\
					\text{③} \quad y = & \\
				\end{array}
				\right.
				\)\\
				\to $y$軸方向の運動は\underline{\hspace{3.5cm}}と一致！
			\end{itemize}
			　\\
			\begin{itembox}[l]{<参考>三角比}
				\begin{minipage}[t]{0.4\linewidth}
					\scalebox{0.8}{\righttriangle{c}{b}{a}{\theta}}
				\end{minipage}
				\begin{minipage}[t]{0.6\linewidth}
					\begin{tabular}{|c|c|c|c|}\hline
						$\theta$ & $\sin\theta$ & $\cos\theta$ & $\tan\theta$ \\\hline
						0° & 0 & 1 & 0 \\\hline
						30° & $\frac{1}{2}$ & $\frac{\sqrt{3}}{2}$ & $\frac{1}{\sqrt{3}}$ \\\hline
						45° & $\frac{\sqrt{2}}{2}$ & $\frac{\sqrt{2}}{2}$ & 1 \\\hline
						60° & $\frac{\sqrt{3}}{2}$ & $\frac{1}{2}$ & $\sqrt{3}$ \\\hline
						90° & 1 & 0 &  \\\hline
					\end{tabular}
					　\\\\\\
				\end{minipage}
			\end{itembox}
		\end{minipage}
	\clearpage

	%=====================================================================================
	\section{演習問題(斜方投射)}
	問1.小球を水平な地面から仰角$30°$、速さ$19.6\,\si{\metre\per\second}$で投げた。
	重力加速度の大きさを$9.8\,\si{\metre\per\second\squared}$,水平方向を　　\,\,$x$軸の正、
	鉛直上向きを$y$軸の正として以下の問いに答えよ。\\
	(1)初速度の$x$成分及び$y$成分は何$\si{\metre\per\second}$か求めよ。ただし、ルートは開かなくてよい。\\\\\\\\\\\\
	\begin{minipage}[t]{0.44\linewidth}
	(2)最高点に達するまでの時間は何$\si{\second}$か求めよ。\\\\\\\\\\\\\\\\\\\\\\
	(3)最高点の高さは何$\si{\metre}$か求めよ。\\\\\\\\\\\\\\\\\\
	\end{minipage}
	\hfill
	\begin{minipage}[t]{0.44\linewidth}
		\vspace{0.5cm}
		\scalebox{0.9}{\simpleobliqueprojectilediagram{19.6\,\mathrm{m/s}}{30}{v_x}{v_y}}
	\end{minipage}\\
	(4)投げてから地面に達するまでの時間は何$\si{\second}$か求めよ。\\\\\\\\\\\\\\\\\\
	(5)小球が地面に落ちたときの水平距離は何$\si{\metre}$か求めよ。ただし、$\sqrt{3} \simeq 1.73$とする。
	% \begin{minipage}[t]{0.44\linewidth}

	% \end{minipage}
	% \hfill
	% \begin{minipage}[t]{0.44\linewidth}

	% \end{minipage}

	\clearpage
\end{document}